\section{The \texorpdfstring{\SxC}{SxC} Decomposition, Confound-Clean Measurement, and a Formal Statement}
\label{sec:method}

\subsection{Setup}
ROME inserts a key$\rightarrow$value association by a rank-one
update to a single MLP down-projection $W$. Given an edit key $k$ (the
subject's last-token key at the critical layer) and target value $v$, the
update is
\begin{equation}
\dW = \frac{(v-Wk)\,k^{\top}}{k^{\top}k}
    = \frac{\mathrm{outer}(v-Wk,\,k)}{\|k\|^2},
\label{eq:rome}
\end{equation}
the minimum-norm solution to $(W+\dW)k=v$.

\subsubsection{Why this is the minimum-norm rank-one solution}
Writing any rank-one update as $\dW = u\,k^{\top}$, the edit constraint
$(W+\dW)k = v$ forces $u = (v-Wk)/(k^{\top}k)$ uniquely. This is also the
minimum-norm solution among \emph{all} matrices $M$ (not only rank-one ones)
with $Mk = v-Wk$: any $N$ with $Nk=0$ adds only orthogonal Frobenius norm, so
the norm is minimized at $N=0$.
We study the identity-covariance (unwhitened) rank-one update. Canonical
ROME's covariance-weighted solve selects a \emph{different} member of the same
solution set $\{M : Mk = v-Wk\}$, replacing $\cos(k,k_p)$ with
$\cos(C^{-1}k,k_p)$; the decomposition of Eq.~\ref{eq:sc} is stated for the
unwhitened update used throughout, and we do not claim it is invariant to
that reweighting.

\subsection{Collateral perturbation to an unrelated probe}
For an unrelated
fact whose key at the same layer is $k_p$, the edit perturbs its value
read-out by $\dW\cdot k_p = (v-Wk)\,(k^{\top}k_p)/\|k\|^2$. Taking magnitudes
and writing $k^{\top}k_p = \|k\|\,\|k_p\|\cos(k,k_p)$,
\begin{equation}
\|\dW\cdot k_p\| = \underbrace{\frac{\|v-Wk\|}{\|k\|}}_{S}\;\|k_p\|\;
                   \underbrace{|\cos(k,k_p)|}_{|C|},
\label{eq:sc}
\end{equation}
where $S$ is a per-\emph{edit} strength (constant across probes), $|C|$ is the
per-(edit,probe) key-geometry, and $\|k_p\|$ is a per-\emph{probe} scalar. The
collateral perturbation factorizes multiplicatively into edit strength, probe
norm, and the edit$\times$probe key-cosine.

\subsubsection{From perturbed read-out to observed damage}
Eq.~\ref{eq:sc} bounds a perturbation to a hidden read-out, not the observed
\emph{damage} we measure -- the drop in the probe's correct-token logit after
the edit. The two are related by the output head, a linear (or near-linear)
map from read-out to logits. Writing $\ell_p$ for the correct-token logit as a
function of the probe's residual-stream contribution $r_p$, a first-order
expansion around the pre-edit read-out gives
\[
\Delta \ell_p \;\approx\; \nabla_{r_p}\ell_p \cdot (\dW\cdot k_p),
\]
so to first order the logit change is a linear functional of $\dW\cdot k_p$,
not merely its norm. Treating $\|\dW\cdot k_p\|$ as informative requires not an
assumption that $\nabla_{r_p}\ell_p$ and $\dW\cdot k_p$ never decouple, but the
weaker, testable claim that the projection factor varies slowly enough across
probes sharing an edit that the \emph{rank order} of $|\Delta\ell_p|$ tracks
that of $\|\dW\cdot k_p\|$, and hence of $S\cdot\|k_p\|\cdot|C|$. This is why we
measure the geometry--damage relationship as a \emph{within-probe Spearman
correlation} -- robust to any monotone map from perturbation magnitude to
logit damage -- rather than asserting exact proportionality between
$\|\dW\cdot k_p\|$ and $\Delta\ell_p$.

\subsection{\texorpdfstring{\SxC}{SxC} is the closed-form reduction of
first-order gradient influence}
Up to the per-probe norm, $S\cdot|C|$ is the exact rank-one reduction of a
first-order (GradSim/TracIn-style) training-influence measure
\citep{pruthi2020gradsim} of an edit on a probe -- obtainable algebraically,
with no backprop; the reduction and the assumptions under which it holds are
made precise in the formal statement of \S\ref{sec:formal} below. At L8 the within-probe $\rho$ for \SxC{} is
\scRhoLeight{}, at L10 \scRhoLten{}, at L12 \scRhoLtwelve{}, and at L14
\scRhoLfourteen{} -- versus raw key-cosine \gWithinLeight{}, \gWithinLten{},
\gWithinLtwelve{}, \gWithinLfourteen{}, respectively. \emph{\SxC{} beats raw
key-cosine at all four layers} once $S$ is normalized per Eq.~\ref{eq:sc}
($S=\|v-Wk\|/\|k\|$, not the unnormalized residual norm alone); the margin is
largest where $S$ carries additional variance (L12/L14). We frame \SxC{} as
the mechanistic, zero-cost surrogate that explains why cosine works and
predicts where its predictive power changes. A closed-form ``GradSim''
baseline built from the same factorization used to define \SxC{} agrees with
\SxC{} by construction, so that agreement is a restatement of the same algebra
rather than an independent empirical check.

To test the surrogate against a genuinely independent quantity, we computed the
true first-order training influence of each edit on each probe with an explicit
backward pass at four layers and three seeds. Three facts hold.
First, the exact read-out identity underlying Eq.~\ref{eq:sc} is verified
programmatically to within numerical tolerance.
% GRADSIM_TRUE_Llama-3.2-1B_L12_s0.json::within_probe_rho.direct_vs_damage.mean = 0.4738 (L12); L8 0.196/0.208/0.212; L10-L14 band 0.43-0.53
Second, this true influence does predict collateral damage: its within-probe
Spearman coupling with damage rises with depth, from $\rho\approx0.20$ at L8 to
$0.43$--$0.53$ across L10--L14 ($0.47$ at the L12 peak).
% GRADSIM_TRUE_Llama-3.2-1B_L12_s0.json::rank_agreement.direct_vs_SC.mean = 0.0874
Third -- the honest limitation -- the within-probe rank agreement between the
true influence and \SxC{} is only about $0.09$, so \SxC{} is \emph{not} a
faithful rank-surrogate of the true first-order influence. We therefore present
\SxC{} not as a stand-in that reproduces influence rankings, but as a zero-cost
\emph{rank estimator} whose value is its closed form: it is exact algebra
(Eq.~\ref{eq:sc}), it explains why key-cosine tracks damage, and it costs no
gradient, while a faithful reconstruction of true-influence rankings remains
beyond its reach.

\subsection{Toward a formal statement}
\label{sec:formal}
The statement below is deliberately conservative: an exact algebraic
reduction at the read-out level, and only a conditional rank-estimator
relationship to observed damage. The error term for ROME's iterative value
solve is carried inside the proof; a numerical bound for production edits is
left to future work.

\iftnnls
\begin{proposition}[\SxC{} as first-order training influence]
\label{prop:sxc-gradsim}
Let the edit at layer $L$ be realized by the rank-one update $\dW = rk^{\top}/(k^{\top}k)$
of Eq.~\ref{eq:rome} for some value $v$ (residual $r=v-Wk$), and let $\ell_p$ denote a
probe's correct-token logit, assumed (Assumption~A2) to depend on $W$ only through the
probe's own read-out $r_p=Wk_p$. Define the first-order training influence of the edit on
the probe as $\mathrm{infl}(e,p):=\langle\nabla_W\ell_p,\dW\rangle_F$, with $\nabla_W\ell_p$
evaluated at the pre-edit weights (Assumption~A3, shared across the independently-applied
edit population). Then, exactly (given A1--A3, no approximation),
\[
\mathrm{infl}(e,p) \;=\; \alpha(e,p)\cdot S\cdot\|k_p\|\cdot C,
\]
where $\alpha(e,p) := \cos(g_p, r)\,\|g_p\|$ with $g_p:=\nabla_{r_p}\ell_p$,
and $S=\|r\|/\|k\|$ and $C=\cos(k,k_p)$ are exactly the terms of the closed-form
decomposition (Eq.~\ref{eq:sc}). The coefficient $\alpha(e,p)$ is fixed in its dependence
on $g_p$ for a given probe $p$ (Assumption~A3) but varies with edit $e$ through
$\cos(g_p,r_e)$, so $\mathrm{infl}(e,p)$ is proportional to, not numerically identical
with, $S_e\|k_p\|C_{e,p}$ as $e$ ranges over the edit population.
In the measured cells the key cosines are effectively non-negative
(negative-cosine key mass is negligible), so $C = |C|$ there and the signed
identity may be read against the reported $|C|$ statistic; the identity
itself is signed. A per-edit-varying
signed scalar does not by itself preserve rank order, so exact coincidence of the
within-probe rank orderings of $\mathrm{infl}(\cdot,p)$ and $S\cdot|C|$ additionally
requires (Assumption~A4$'$) that $\alpha(e,p)$ is sign-consistent across edits for fixed
$p$ \emph{and} does not covary with $S_e|C_{e,p}|$ enough to reorder pairs -- a
plausibility condition, not a consequence of A1--A3 (and one our cells violate:
per-probe sign-consistency of $\alpha$ is $0.66$--$0.81$ against a $0.50$ chance
floor, with median coefficient of variation $1.0$--$2.4$; see the remark below). Under A1--A3 and A4$'$, the
within-probe Spearman coupling this paper reports between $|C|$ and collateral damage is
a rank \emph{estimator} of the within-probe Spearman coupling between first-order
training influence and damage -- not an exact identity of the two statistics.
Since A4$'$ does not hold in our cells, the conditional's antecedent fails and we
do not invoke it: the proposition's standing in this paper is the exact
read-out identity of Eq.~\ref{eq:sc} alone, which needs A1--A3 only.
\end{proposition}
\begin{proof}[Proof sketch]
Substitute $\dW$ and $\nabla_W\ell_p=g_pk_p^{\top}$ (chain rule under A2) into the
Frobenius inner product, using $\langle ab^\top,cd^\top\rangle_F=(a^\top c)(b^\top d)$ and
$k^\top k_p=\|k\|\|k_p\|C$. Under A4$'$, a per-column ($p$ fixed),
sign-consistent multiplier would leave Spearman rank correlation
approximately unchanged (exactly unchanged only if $\alpha(e,p)$ were constant, which is
not assumed) -- hence ``estimator,'' not ``coincides.'' We do not quantify the
resulting rank perturbation; the step is heuristic, and the measurement reported
below shows the multiplier's variation is in fact large enough to destroy the
agreement. An explicit error term for the
additional deviation between this proposition's single-step realization and the
iterative multi-step value solve used throughout the empirical sections follows the
\emph{same} substitution with $\Delta v:=v_{\text{iter}}-v_{\text{conv}}$ in place of $r$,
where $v_{\text{iter}}$ is the value returned by the regularized multi-step solve and
$v_{\text{conv}}$ the fully converged, unregularized solve:
\begin{align*}
&\mathrm{infl}_{\text{iter}}(e,p)-\mathrm{infl}_{\text{conv}}(e,p)\\
&\quad= \big[\cos(g_p,\Delta v)\,\|g_p\|\big]\cdot \Delta S\cdot\|k_p\|\cdot C,
\end{align*}
with $\Delta S := \|\Delta v\|/\|k\|$,
since $\dW_{\text{iter}}-\dW_{\text{conv}}=\Delta v\,k^{\top}/\|k\|^2$ has the same
rank-one form as $\dW$. Taking magnitudes and using $|\cos(g_p,\Delta v)|\le 1$ gives
$\big|\mathrm{infl}_{\text{iter}}-\mathrm{infl}_{\text{conv}}\big|\le \|g_p\|\,\Delta S\,\|k_p\|\,|C|$,
exact given the definition of $\mathrm{infl}$ (no further linearization beyond A1--A3).
\end{proof}
\fi
\iftaslp
\begin{remark}[\SxC{} and first-order training influence]
The closed-form product $S\cdot\|k_p\|\cdot|C|$ (Eq.~\ref{eq:sc}) is exactly the
perturbation magnitude the edit induces on a probe's pre-logit read-out, and it is
\emph{exactly} proportional -- not merely correlated -- to a TracIn/GradSim-style
first-order training influence of the edit on the probe's loss,
$\langle\nabla_W\ell_p,\dW\rangle_F$, with a proportionality factor that depends on the
alignment between the probe's gradient and the edit's target-residual direction and so
varies across edits for a fixed probe. Treating the within-probe rank correlation this
paper reports between $|C|$ and collateral damage as a rank correlation between
first-order training influence and damage additionally requires (Assumption~A4$'$) that
this factor be sign-consistent across edits for a fixed probe \emph{and} not
systematically reorder edits by their $S\times C$ rank. We test this condition
rather than assume it, and it does not hold in our cells: the per-probe
sign-consistency of $\alpha(e,p)$ is $0.66$--$0.81$ (chance $0.50$) and its
median coefficient of variation is $1.0$--$2.4$ across four layers and three seeds.
A4$'$ is therefore violated, which is exactly why the measured rank agreement
between true influence and \SxC{} is near zero. Consequently we do not claim
the reported coupling as a rank estimator of the influence--damage coupling;
the closed form's standing rests on the exact read-out identity of
Eq.~\ref{eq:sc}, which is independent of A4$'$. An explicit error term for the gap between this single-step
realization and ROME's actual iterative, regularized multi-step value solve follows the
identical factorization (Eq.~\ref{eq:sc} applied to the value-solve gap $\Delta v$,
bounded via the trivial $|\cos(g_p,\Delta v)|\le 1$) and is exact given the definition of
the influence measure, of the same $S\times C$ functional form, scaled by the norm of the
gap between the iterative and converged value vectors; deriving a numerical value for that
gap on production (non-toy) edits is left to future work.
\end{remark}
\fi

\subsection{Confound-clean measurement}
\label{sec:gate}

\subsubsection{Definition (within-probe partialled coupling)}
The empirical test of Eq.~\ref{eq:sc} is a \emph{within-probe} partialled
Spearman between $|C|$ and damage: fixing a probe $p$ and ranging the edit
index $e$ over all edits, the estimator correlates $\{|C_{e,p}|\}_e$ with
$\{\text{damage}_{e,p}\}_e$ down that single probe column. This differences out
$\|k_p\|$ and the probe's intrinsic susceptibility to any edit (both constant
within a column), so only the edit-specific, per-(edit,probe) key geometry can
explain the correlation. Aggregating across all probes (we report the mean
across probe columns, with per-layer standard deviations over three seeds)
yields the reported within-probe $\rho$.

\subsubsection{Procedure (measurement protocol and controls)}
We run ROME on Llama-3.2-1B \citep{grattafiori2024llama3},
CounterFact, $N{=}200$ edits
$\times$ $M{=}500$ probes, with two confound controls: the \emph{known-fact
filter} (restrict to probes the base model assigns the true answer,
$\text{pre\_p}>0.05$) and the \emph{successful-edit filter} (drop failed edits).
This retires two objections the design pre-registered as
mandatory. First, \emph{probe-marginal leakage}: if the within-probe estimator
were simply recovering ``which probes are fragile regardless of the edit,''
the pooled correlation (which does not condition on probe identity) should be
far larger than the within-probe one. In fact the pooled $\rho$ ($\sim$0.40 at
L8) barely exceeds the within-probe $\rho$ (\gWithinLeight{}), so the signal
is edit-specific pairwise geometry, not probe-marginal fragility. Second,
\emph{non-independence} across probes sharing the same edit set: a
column-permutation null (independently shuffling damage values within each
probe column, which destroys the edit-to-probe pairing while preserving each
probe's marginal damage distribution) gives a permutation-test $p$-value at
the \gPermFloor{} floor, i.e.\ the true pairing is essentially never matched
by chance under this null. The signal further survives row-marginal
double-centering (subtracting both the per-edit and per-probe mean damage
before correlating) and survives partialling matrix-norm-growth out of each
column, i.e.\ it is not an artifact of a confound shared between key-cosine
and the norm-growth quantity examined in Section~\ref{sec:regime}. Full gate
results are in \S\ref{sec:regime} and Figure~\ref{fig:layerlaw}.
